%% Question bank chapter file - generated by bank_tools.py build (do not edit while the bank is still being built:
%% it is regenerated from the per-test data in Question Bank\data). Every question is one bankq block:
%%   \begin{bankq}[id=..., source={...}, paper=.., qno=.., type=..., subtopic={...}, answer={...}, answer_src=key|solved,
%%                 working={...}, difficulty=1-10, diff_parts={K S I C F}, diff_note={...},
%%                 group=<passage id or empty>, also_in={duplicate copies in other tests}]
%% difficulty = min(10, K+S+I+C+F-1): K concepts 1-3, S steps 1-4, I insight 0-2, C computation 0-2, F format 0-1
%%   ...question LaTeX, options, figures (\bfig / \bfigi -> img/)...
%%   \end{bankq}
%% Passages shared by several questions are bankpassage blocks; their questions carry group=<passage id>.
\documentclass[10pt]{article}
\usepackage[a4paper,margin=1.8cm]{geometry}
\usepackage[T1]{fontenc}
\usepackage{amsmath,amssymb,mathtools}
\usepackage{graphicx}
\usepackage[export]{adjustbox}
\usepackage{xcolor,needspace,keyval}
\usepackage{enumitem}
\usepackage{booktabs,array,tabularx}
\usepackage[hidelinks]{hyperref}
\setlength{\parindent}{0pt}
\setlength{\parskip}{3pt}
\makeatletter
\define@key{bq}{id}{\def\bq@id{#1}}\define@key{bq}{source}{\def\bq@src{#1}}\define@key{bq}{paper}{\def\bq@paper{#1}}
\define@key{bq}{qno}{\def\bq@qno{#1}}\define@key{bq}{type}{\def\bq@type{#1}}\define@key{bq}{subtopic}{\def\bq@sub{#1}}
\define@key{bq}{answer}{\def\bq@ans{#1}}\define@key{bq}{answer_src}{\def\bq@asrc{#1}}\define@key{bq}{working}{\def\bq@work{#1}}
\define@key{bq}{group}{\def\bq@group{#1}}\define@key{bq}{also_in}{\def\bq@also{#1}}
\define@key{bq}{difficulty}{\def\bq@diff{#1}}\define@key{bq}{diff_parts}{\def\bq@dparts{#1}}\define@key{bq}{diff_note}{\def\bq@dnote{#1}}
\newenvironment{bankq}[1][]{\def\bq@work{}\def\bq@also{}\def\bq@group{}\def\bq@diff{?}\def\bq@dparts{}\def\bq@dnote{}\setkeys{bq}{#1}%
  \par\needspace{4\baselineskip}\medskip\noindent{\small\textbf{\bq@id}\hfill\color{gray}[\bq@src, Paper \bq@paper, Q\bq@qno]
  \bq@type\ $\cdot$ \bq@sub\ $\cdot$ \textbf{D\bq@diff}}\par\nopagebreak\relax}
  {\par{\small\color{teal}Answer: \textbf{\bq@ans}\ifx\bq@work\@empty\else\ (solved: \bq@work)\fi
  \ifx\bq@also\@empty\else\par\color{gray}Also in: \bq@also\fi}\par}
\define@key{bp}{id}{\def\bp@id{#1}}\define@key{bp}{source}{\def\bp@src{#1}}\define@key{bp}{questions}{\def\bp@qs{#1}}
\newenvironment{bankpassage}[1][]{\setkeys{bp}{#1}\par\needspace{4\baselineskip}\bigskip\noindent
  {\small\textbf{Passage \bp@id}\hfill\color{gray}[\bp@src; questions \bp@qs]}\par\nopagebreak\relax\itshape}{\par}
\makeatother
\newcommand{\bfig}[2]{\par\smallskip{\centering\includegraphics[width=#1pt,max width=0.62\linewidth,max height=0.32\textheight]{img/#2}\par}\smallskip}
\newcommand{\bfigi}[2]{\includegraphics[width=#1pt,max width=\linewidth,max height=0.32\textheight,valign=c]{img/#2}}
\newcommand{\bmissing}[1]{\textit{[figure missing in the source conversion: \detokenize{#1}]}}
\newenvironment{qoptions}{\begin{itemize}[leftmargin=1.8em,itemsep=1pt,topsep=2pt,label={}]}{\end{itemize}}
\newcommand{\opt}[2]{\item[(#1)] #2}
\providecommand{\fwstart}{}\providecommand{\fwend}{}\providecommand{\qfullsplit}{}
\newcolumntype{L}[1]{>{\raggedright\arraybackslash}p{#1}}
\newcolumntype{C}[1]{>{\centering\arraybackslash}p{#1}}
\newcolumntype{R}[1]{>{\raggedleft\arraybackslash}p{#1}}
\providecommand{\sectionhdr}[1]{}\providecommand{\subsectionhdr}[1]{}


\begin{document}
{\Large\bfseries PnC --- Question Bank}\par\smallskip
{\small Mathematics chapter M04. Questions from all converted test papers, grouped by source series. Answers: printed key where it matches the paper, otherwise solved (answer\_src=solved, with working). Exact repeats across papers appear once, with also\_in listing the other copies.}\par\medskip\hrule

\section*{Nurture SRG papers}
\begin{bankq}[id=NSRG02-P1-MATH-S1-Q01, source={Nurture SRG - Test 02 (06 Dec 2023)}, paper=1, qno=1, type=single, subtopic={Permutations}, answer={B}, answer_src=solved, working={Colourings of a 4-cycle with 6 values: $(6-1)^4+(6-1)=630$}, difficulty=6, diff_parts={K1 S2 I2 C2 F0}, diff_note={concepts: colourings of a cycle with a dice | steps: cycle of four; adjacent differ; counting formula; value}]
4 people sit around a circular table and everyone rolls a dice. The number of ways in which they can roll the dice such that no two people next to each other roll the same number is
\begin{qoptions}
\opt{A}{750}
\opt{B}{630}
\opt{C}{1296}
\opt{D}{None of these}
\end{qoptions}
\end{bankq}
\begin{bankq}[id=NSRG02-P1-MATH-S1-Q04, source={Nurture SRG - Test 02 (06 Dec 2023)}, paper=1, qno=4, type=single, subtopic={Combinations}, answer={C}, answer_src=solved, working={Increasing five digit numbers are $\binom95=126$; 70 start with 1, 20 more with 23, then 6 with 245, so the 97th is the first number of the form 246xx, which is 24678}, difficulty=5, diff_parts={K1 S2 I2 C1 F0}, diff_note={concepts: rank of a number with increasing digits | steps: count by leading digits; narrow down; find the number}]
All the five digit numbers in which each successive digit exceeds its predecessor are arranged in ascending order. The 97th number is
\begin{qoptions}
\opt{A}{24589}
\opt{B}{24679}
\opt{C}{24678}
\opt{D}{24689}
\end{qoptions}
\end{bankq}
\begin{bankq}[id=NSRG02-P1-MATH-S2-Q05, source={Nurture SRG - Test 02 (06 Dec 2023)}, paper=1, qno=5, type=multi, subtopic={Distributions}, answer={B,C,D}, answer_src=solved, working={Distributions with at most two per room: $2220=2^2\cdot3\cdot5\cdot37$; digit sum 6, not divisible by 7, 24 factors; not a square}, difficulty=7, diff_parts={K1 S2 I2 C2 F1}, diff_note={concepts: distribution of friends into rooms with limits | steps: cases by occupancy; arrangements of each case; total; properties of N}]
5 friends can be accommodated in 5 different rooms of a hotel in $N$ ways such that there are not more than 2 friends in any of the rooms, where
\begin{qoptions}
\opt{A}{$N$ is a perfect square}
\opt{B}{Sum of digits of $N$ is 6}
\opt{C}{$N$ is not divisible by 7}
\opt{D}{$N$ has 24 factors}
\end{qoptions}
\end{bankq}
\begin{bankq}[id=NSRG02-P1-MATH-S2-Q06, source={Nurture SRG - Test 02 (06 Dec 2023)}, paper=1, qno=6, type=multi, subtopic={Combinations}, answer={B,C,D}, answer_src=solved, working={Counting triplets $a<b<c\le15$ with common factor: $N=46$, digit sum 10, largest single digit factor is 2}, difficulty=7, diff_parts={K1 S2 I2 C2 F1}, diff_note={concepts: sorted triplets with a common factor | steps: count by the common divisor; inclusion exclusion; total; properties}]
We say that a triplet of integers $(a,b,c)$ is sorted if $a<b<c$. Suppose there are $N$ number of sorted triplets of positive integers such that $c\le15$ and the greatest common divisor of $a,b,c$ is greater than 1. Then which of the following is INCORRECT?
\begin{qoptions}
\opt{A}{Sum of digits of $N$ is 10}
\opt{B}{Sum of the digits of $N$ is 9}
\opt{C}{Largest single digit factor of $N$ is 6}
\opt{D}{Largest single digit factor of $N$ is 7}
\end{qoptions}
\end{bankq}
\begin{bankq}[id=NSRG02-P1-MATH-S2-Q08, source={Nurture SRG - Test 02 (06 Dec 2023)}, paper=1, qno=8, type=multi, subtopic={Counting principles}, answer={A,B,D}, answer_src=solved, working={$m=7^9$ (10 factors); $n=7^9-2\cdot6^9+5^9=22151340$, $m-n$ odd, $n$ divisible by 6 but not a factor of $m$}, difficulty=7, diff_parts={K1 S2 I2 C2 F1}, diff_note={concepts: counting ordered triples of subsets with conditions | steps: choices per element; inclusion exclusion for the conditions; values of m and n; check properties}]
A, B, C are subsets of $\{1,2,3,\dots,9\}$. The number of ordered triplets $(A,B,C)$ such that $A\cap B\cap C=\phi$ is $m$ and number of ordered triplets $(A,B,C)$ such that $A\cap B\cap C=\phi$, $A\cap B\ne\phi$ and $B\cap C\ne\phi$ is $n$, then
\begin{qoptions}
\opt{A}{Number of factors of $m$ is 10}
\opt{B}{$m-n$ is odd number}
\opt{C}{$n$ is a factor of $m$}
\opt{D}{$n$ is a multiple of 6}
\end{qoptions}
\end{bankq}
\begin{bankq}[id=NSRG02-P1-MATH-S2-Q09, source={Nurture SRG - Test 02 (06 Dec 2023)}, paper=1, qno=9, type=multi, subtopic={Combinations}, answer={A,B,D}, answer_src=solved, working={Arrangements with no 7 consecutive whites $=\binom{28}{14}-15\binom{21}7+\binom{15}{2}$, so $k=14$, $\ell=15$, $m=2$ or 13}, difficulty=8, diff_parts={K1 S3 I2 C2 F1}, diff_note={concepts: arrangements of balls with no long run | steps: total arrangements; complement; counting runs of whites; formula; match parameters}]
14 identical black balls and 14 identical white balls are to be arranged in a row. The number of ways to arrange them such that on picking any 7 consecutive balls we get atleast one black ball is ${}^{28}c_k-\ell\cdot{}^{21}c_7+{}^{15}c_m$, then
\begin{qoptions}
\opt{A}{$k+\ell$ can be 29}
\opt{B}{$k-m$ can be 12}
\opt{C}{$k=19$}
\opt{D}{$m,k,\ell$ can be three consecutive term of an A.P.}
\end{qoptions}
\end{bankq}
\begin{bankq}[id=NSRG02-P1-MATH-S2-Q10, source={Nurture SRG - Test 02 (06 Dec 2023)}, paper=1, qno=10, type=multi, subtopic={Counting principles}, answer={B,C,D}, answer_src=solved, working={$a_n=a_{n-1}+n$, $a_n=1+\tfrac{n(n+1)}2$, $a_6=22$, $a_{10}=56$}, difficulty=6, diff_parts={K1 S2 I2 C1 F1}, diff_note={concepts: regions formed by lines in a plane | steps: recurrence for regions; values; check}]
There are $n$ lines in a plane, no two of which are parallel and no 3 of them are concurrent. Let the plane be divided by $n$ lines in $a_n$ parts, then
\begin{qoptions}
\opt{A}{$a_n=a_{n-1}+(n-1)$}
\opt{B}{$a_n=a_{n-1}+n$}
\opt{C}{$a_6=22$}
\opt{D}{$a_{10}=56$}
\end{qoptions}
\end{bankq}
\begin{bankq}[id=NSRG02-P1-MATH-S2-Q11, source={Nurture SRG - Test 02 (06 Dec 2023)}, paper=1, qno=11, type=multi, subtopic={Combinations}, answer={B,D}, answer_src=solved, working={Selections equal the coefficient of $x^4$ in $(1+x)^3(1+x+x^2)^2$, which is 18}, difficulty=7, diff_parts={K1 S2 I2 C2 F1}, diff_note={concepts: selections of letters from a word with repetition | steps: letter multiplicities; generating function; coefficient; compare}]
The number of selections of 4 letters from the word COLLEGE is equal to
\begin{qoptions}
\opt{A}{24}
\opt{B}{18}
\opt{C}{Coefficient of $x^4$ in $(1+x)^2\left(1+x+x^2\right)^3$}
\opt{D}{Coefficient of $x^4$ in $(1+x)^3\left(1+x+x^2\right)^2$}
\end{qoptions}
\end{bankq}
\begin{bankq}[id=NSRG02-P1-MATH-S2-Q12, source={Nurture SRG - Test 02 (06 Dec 2023)}, paper=1, qno=12, type=multi, subtopic={Derangements}, answer={A,D}, answer_src=solved, working={$N=120-11=109$ (prime, digit sum 10); the harmonic mean of the divisors of 112 is not 5}, difficulty=7, diff_parts={K1 S2 I2 C2 F1}, diff_note={concepts: permutations with many displaced positions | steps: total permutations; those with three or more fixed; count; properties of the number}]
Let $N$ be the number of arrangements of all the digits of the number 12345 such that atleast 3 digits will not come in it's original position, then
\begin{qoptions}
\opt{A}{Sum of digits of $N$ is 10}
\opt{B}{Number of ways to resolve $N$ as product of two nos. is 6}
\opt{C}{Harmonic mean of all divisors of $(N+3)$ is 5}
\opt{D}{$N$ is a prime number}
\end{qoptions}
\end{bankq}
\begin{bankq}[id=NSRG02-P1-MATH-S3-Q01, source={Nurture SRG - Test 02 (06 Dec 2023)}, paper=1, qno=1, type=numerical, subtopic={Combinations}, answer={9}, answer_src=solved, working={$N=\binom{15}7=6435=3^2\cdot5\cdot11\cdot13$, largest single digit factor 9}, difficulty=6, diff_parts={K1 S2 I2 C1 F1}, diff_note={concepts: seven digit numbers with digits in order | steps: combinations with repetition; value; factor}]
Let $N$ denote the number of 7 digit positive integers having the property that their digits are in non decreasing order. Find the largest single digit factor of $N$
\end{bankq}
\begin{bankq}[id=NSRG02-P1-MATH-S3-Q04, source={Nurture SRG - Test 02 (06 Dec 2023)}, paper=1, qno=4, type=numerical, subtopic={Counting principles}, answer={8}, answer_src=solved, working={Each letter after the first has 2 allowed choices: $3\cdot2^{n-1}=384$, so $n=8$}, difficulty=6, diff_parts={K1 S2 I2 C1 F1}, diff_note={concepts: words with forbidden pairs | steps: choices at each step; formula; solve for n}]
The number of $n$ - letter words consisting only of the letters A, B and C such that A is never immediately followed by B, B is never immediately followed by C and C is never immediately followed by A are 384. Find value of $n$.
\end{bankq}
\begin{bankq}[id=NSRG02-P1-MATH-S3-Q06, source={Nurture SRG - Test 02 (06 Dec 2023)}, paper=1, qno=6, type=numerical, subtopic={Counting principles}, answer={1960}, answer_src=solved, working={Digits sum to 73 so the deficiency from nine nines is 8; counting those with alternating sum divisible by 11 gives 1960}, difficulty=7, diff_parts={K1 S2 I2 C2 F1}, diff_note={concepts: nine digit numbers with given sum and divisibility by eleven | steps: deficiency from all nines; alternating sum condition; enumerate the cases; count}]
Find the number of nine digit numbers divisible by 11 and having sum of digits as 73.
\end{bankq}
\begin{bankq}[id=NSRG03-P1-MATH-S2-Q10, source={Nurture SRG - Test 03 (04 Feb 2024)}, paper=1, qno=10, type=multi, subtopic={Counting principles}, answer={A,B}, answer_src=solved, working={Missing pair $(16,17)$: $N=2^3\cdot3^2\cdot5\cdot7\cdot11\cdot13\cdot19$, $K=384$}, difficulty=8, diff_parts={K1 S3 I2 C2 F1}, diff_note={concepts: smallest integer divisible by all but two consecutive numbers | steps: which pair can be missing; smallest number; factorisation; number of divisors; compare}]
Let $S=\{1,2,3,\dots,20\}$ be the set of all positive integers from 1 to 20. Suppose that $N$ is the smallest positive integer such that exactly eighteen numbers from $S$ are factors of $N$, and the only two numbers from $S$ that are not factors of $N$ are consecutive integers. If $K$ is the total number of positive integral divisors of $N$, then $K$ is greater than
\begin{qoptions}
\opt{A}{200}
\opt{B}{300}
\opt{C}{400}
\opt{D}{500}
\end{qoptions}
\end{bankq}
\begin{bankq}[id=NSRG03-P1-MATH-S3-Q01, source={Nurture SRG - Test 03 (04 Feb 2024)}, paper=1, qno=1, type=numerical, subtopic={Distributions}, answer={5246}, answer_src=solved, working={$p=\binom84\binom84=4900$; $q=\sum\binom4a^3=346$; $p+q=5246$}, difficulty=7, diff_parts={K1 S2 I2 C2 F1}, diff_note={concepts: words with restrictions on positions of letters | steps: count when no X in the first block; case counting for the second condition; add}]
Suppose we are constructing 12 lettered word with four X's four Y's and four Z's. Let $p$ denotes the number of 12 lettered words which have no X's in the first four letters, and $q$ denotes the number of 12 lettered words which have no X's in the first four letters, no Y's in the middle four letters and no Z's in the last four letters. Find $(p+q)$.
\end{bankq}
\begin{bankq}[id=NSRG03-P1-MATH-S3-Q04, source={Nurture SRG - Test 03 (04 Feb 2024)}, paper=1, qno=4, type=numerical, subtopic={Distributions}, answer={4}, answer_src=solved, working={Number of ways $=826$, $A=8,B=2,C=6$, so $A+B-C=4$}, difficulty=8, diff_parts={K1 S3 I2 C2 F1}, diff_note={concepts: distributions of marks with bounds | steps: shift variables; inclusion exclusion; count; read the digits; combine}]
In a paper of mathematics there are 5 questions, such that the sum of the marks is 30 and the marks for any question is not less than 2 and not more than 8 (only integer marks are awarded). If the number of ways in which the marks can be awarded is $ABC$, where $A,B,C\in N$ and $A,B,C\le9$. Then value of $A+B-C$ is \underline{\hspace{1cm}}
\end{bankq}
\begin{bankq}[id=NSRG03-P2-MATH-S1-Q02, source={Nurture SRG - Test 03 (04 Feb 2024)}, paper=2, qno=2, type=single, subtopic={Combinations}, answer={A}, answer_src=solved, working={Choose 7 of the 10 squares covering all four rows: $120-45+1=76$; arrange the letters $\dfrac{7!}{2!2!}=1260$; $76\times1260=95760$}, difficulty=6, diff_parts={K1 S2 I2 C2 F0}, diff_note={concepts: arranging letters in a staircase of squares with no empty row | steps: choose the squares; count the selections covering all rows; arrange the letters; multiply}]
The number of ways in which all the letters of the word `BIGBOSS' can be placed in the squares of the figure shown, so that no row remains empty is equal to\par
\bfig{73.2}{NSRG03_p2_math_s1_q02_1.png}
\begin{qoptions}
\opt{A}{95760}
\opt{B}{94640}
\opt{C}{97840}
\opt{D}{82120}
\end{qoptions}
\end{bankq}
\begin{bankq}[id=NSRG03-P2-MATH-S2-Q09, source={Nurture SRG - Test 03 (04 Feb 2024)}, paper=2, qno=9, type=multi, subtopic={Combinations}, answer={B,C}, answer_src=solved, working={Counting by computer: even $x$ gives 16065, odd $x$ gives 14722}, difficulty=7, diff_parts={K1 S2 I2 C2 F1}, diff_note={concepts: ordered triples with an inequality and parity | steps: count triples with the bound; separate by parity of x; compute each count}]
The number of ordered triplets $(x,y,z)$ of non negative integers satisfying the conditions. $x+y+z\le100$ and $x\le y\le z$
\begin{qoptions}
\opt{A}{if $x$ is odd is 14724}
\opt{B}{if $x$ is odd is 14722}
\opt{C}{if $x$ is even is 16065}
\opt{D}{if $x$ is even is 16164}
\end{qoptions}
\end{bankq}
\begin{bankq}[id=NSRG03-P2-MATH-S2-Q10, source={Nurture SRG - Test 03 (04 Feb 2024)}, paper=2, qno=10, type=multi, subtopic={Combinations}, answer={A,C,D}, answer_src=solved, working={The 105th number is the last with leading digit 2: 26789}, difficulty=6, diff_parts={K1 S2 I2 C1 F1}, diff_note={concepts: 105th number with increasing digits | steps: count by leading digits; find the number; check the digits}]
All the five digit numbers in which each successive digit exceeds its predecessor are arranged in the increasing order. The 105th number does not contain the digit
\begin{qoptions}
\opt{A}{1}
\opt{B}{2}
\opt{C}{3}
\opt{D}{4}
\end{qoptions}
\end{bankq}
\begin{bankq}[id=NSRG03-P2-MATH-S3-Q02, source={Nurture SRG - Test 03 (04 Feb 2024)}, paper=2, qno=2, type=numerical, subtopic={Permutations}, answer={161}, answer_src=solved, working={Counting words that precede CALCU in dictionary order gives rank 161}, difficulty=7, diff_parts={K1 S2 I2 C2 F1}, diff_note={concepts: rank of a word in dictionary order | steps: letters available; words that start with earlier letters; words with the same start; count}]
Let 5 letters words are formed using the letters of the word `CALCULUS' and $N$ is the rank of word CALCU among them as arranged in dictionary. Then $N$ is
\end{bankq}

\section*{TOAS papers}
\begin{bankq}[id=TOAS01-P1-MATH-S1-Q03, source={TOAS - Test 01 (30 Jan 2022)}, paper=1, qno=3, type=single, subtopic={Distributions}, answer={C}, answer_src=key, difficulty=2, diff_parts={K1 S1 I1 C0 F0}, diff_note={concepts: number of solutions of a linear equation | steps: x = 2a + 1, y = 2b, z = 2c reduces to a + b + c = 23; C(25,2) = 300 | I: substitution}]
If $x$ is an odd integer and $y,z$ are even integers, satisfying $x+y+z=47$ and $x,y,z\ge0$, then number of solutions are
\begin{qoptions}
\opt{A}{100}
\opt{B}{200}
\opt{C}{300}
\opt{D}{400}
\end{qoptions}
\end{bankq}
\begin{bankq}[id=TOAS01-P1-MATH-S2-Q10, source={TOAS - Test 01 (30 Jan 2022)}, paper=1, qno=10, type=multi, subtopic={Counting principles}, answer={A,B,C,D}, answer_src=key, difficulty=3, diff_parts={K1 S2 I0 C0 F1}, diff_note={concepts: Legendre's formula | steps: exponent of 5 in 500! is 124; compare with exponent of 2; evenness}]
If $\underline{|500}=k.\left(2^a3^b5^c\right)$, (where $a,b,c\in N$, $k$ does not contain any factor of 2,3 and 5) then-
\begin{qoptions}
\opt{A}{$a>c$}
\opt{B}{$c=124$}
\opt{C}{number of zeros at the end of $\underline{|500}$ is equal to 124}
\opt{D}{$\underline{|500}$ is an even number.}
\end{qoptions}
\end{bankq}
\begin{bankpassage}[id=TOAS01-P1-MATH-S3-G1-2, source={TOAS - Test 01 (30 Jan 2022)}, questions={1-2}]
Consider the name of a baby girl `HARSHITA'. New names are formed (with or without meaning) by using all the letters of the given word.
\end{bankpassage}
\begin{bankq}[id=TOAS01-P1-MATH-S3-Q01, source={TOAS - Test 01 (30 Jan 2022)}, paper=1, qno=1, type=numerical, subtopic={Permutations}, answer={720}, answer_src=key, difficulty=3, diff_parts={K1 S2 I0 C0 F1}, diff_note={concepts: arrangements with repetition | steps: choose the even places for the vowels; arrange AAI; arrange the consonants with H repeated}, group=TOAS01-P1-MATH-S3-G1-2]
Number of words in which vowels occupy even places -
\end{bankq}
\begin{bankq}[id=TOAS01-P1-MATH-S3-Q02, source={TOAS - Test 01 (30 Jan 2022)}, paper=1, qno=2, type=numerical, subtopic={Permutations}, answer={360}, answer_src=key, difficulty=3, diff_parts={K1 S1 I1 C0 F1}, diff_note={concepts: arrangements with a block | steps: treat AAI as one unit; arrange 6 units with H repeated | I: block method}, group=TOAS01-P1-MATH-S3-G1-2]
Number of words which contains the word `AAI' is-
\end{bankq}
\begin{bankq}[id=TOAS01-P1-MATH-S4-Q07, source={TOAS - Test 01 (30 Jan 2022)}, paper=1, qno=7, type=numerical, subtopic={Counting principles}, answer={5}, answer_src=key, difficulty=4, diff_parts={K2 S2 I0 C0 F1}, diff_note={concepts: binomial coefficient sum, divisor counting | steps: n = 10 from 2\^{}n = 1024; factorise 1200 and impose the factor 75; count the powers of 2}]
Let sum of binomial coefficient in the expansion of $(x+2y)^n$ is 1024 ($n\in N$), then number of divisors of $(120)(n)$ divisible by 75 is
\end{bankq}
\begin{bankq}[id=TOAS01-P1-MATH-S4-Q09, source={TOAS - Test 01 (30 Jan 2022)}, paper=1, qno=9, type=numerical, subtopic={Counting principles}, answer={7}, answer_src=key, difficulty=4, diff_parts={K1 S2 I1 C0 F1}, diff_note={concepts: LCM counting via exponents | steps: exponent triples for the prime 5; exponent triples for the prime 3; multiply the counts | I: prime-wise counting}]
If L.C.M. of $(a,b)=5^2\times3^4$, L.C.M. of $(b,c)=5^3\times3^4$ \& L.C.M. of $(c,a)=5^3\times3^4$, then possible ordered triplets of $(a,b,c)$ are $k$ then $(k-58)$ is equal to
\end{bankq}
\begin{bankq}[id=TOAS01-P2-MATH-S1-Q12, source={TOAS - Test 01 (30 Jan 2022)}, paper=2, qno=12, type=single, subtopic={Permutations}, answer={B}, answer_src=key, difficulty=2, diff_parts={K1 S1 I1 C0 F0}, diff_note={concepts: counting with exclusion | steps: total arrangements minus those with the two papers together; rewrite in the option form | I: complementary counting}]
Number of ways in which 10 examination papers can be arranged so that the best and the worst papers never come together, is -
\begin{qoptions}
\opt{A}{$2({}^9C_2)$}
\opt{B}{$2({}^9C_2)8!$}
\opt{C}{$2({}^9C_2)7!$}
\opt{D}{$8!$}
\end{qoptions}
\end{bankq}
\begin{bankq}[id=TOAS01-P2-MATH-S1-Q13, source={TOAS - Test 01 (30 Jan 2022)}, paper=2, qno=13, type=single, subtopic={Counting principles}, answer={C}, answer_src=key, difficulty=3, diff_parts={K1 S2 I1 C0 F0}, diff_note={concepts: LCM and prime-wise counting | steps: prime factorisation of 150; triples for each prime; product; check the options | I: exponent triples}]
LCM of 3 natural numbers $a,b,c$ is 150. If $N$ is the number of such ordered triplet, then -
\begin{qoptions}
\opt{A}{$N$ is multiple of 5}
\opt{B}{$N$ is an even number}
\opt{C}{$N$ is multiple of 7}
\opt{D}{$N$ is greater then 1000}
\end{qoptions}
\end{bankq}
\begin{bankq}[id=TOAS01-P2-MATH-S2-Q04, source={TOAS - Test 01 (30 Jan 2022)}, paper=2, qno=4, type=numerical, subtopic={Counting principles}, answer={8}, answer_src=solved, working={Ordered compositions of 10 into parts 2,3,4: $f(n)=f(n-2)+f(n-3)+f(n-4)$ gives $f(10)=17$; $[17/2]=8$}, difficulty=4, diff_parts={K1 S2 I1 C0 F1}, diff_note={concepts: compositions of an integer | steps: recurrence f(n) = f(n-2) + f(n-3) + f(n-4); base values; build up to f(10); halve | I: recurrence}]
The postal stamps of three different valuations viz Rs 2, Rs 3 and Rs 4 are available. If the number of ways in which these stamps can be posted on envelop one after the other so that total valuation of stamps on envelop is Rs 10 is $N$, then $\left[\dfrac N2\right]$ is equal to (where [.] denotes greatest integer function)
\end{bankq}
\begin{bankq}[id=TOAS01-P2-MATH-S2-Q05, source={TOAS - Test 01 (30 Jan 2022)}, paper=2, qno=5, type=numerical, subtopic={Permutations}, answer={0}, answer_src=solved, working={Middle digit $m$ a square: $m=4$: $\binom33\binom43=4$; $m=9$: $\binom83\binom93=56\cdot84=4704$ (first digit nonzero); $\lambda=4708$, tens digit 0}, difficulty=5, diff_parts={K1 S3 I1 C0 F1}, diff_note={concepts: digit counting with ordering constraints | steps: middle digit options; ascending choices; descending choices; handle leading digit; add | I: ordered choices are combinations}]
If number of seven digit numbers in which first four digits are in ascending order, last four digits are in descending order and middle digit is a perfect square is $\lambda$, then digit at the tens place of $\lambda$ is
\end{bankq}
\begin{bankq}[id=TOAS01-P2-MATH-S2-Q06, source={TOAS - Test 01 (30 Jan 2022)}, paper=2, qno=6, type=numerical, subtopic={Counting principles}, answer={4}, answer_src=solved, working={Numbers: 11,12,15,22,24,33,36,44,48,55,66,77,88,99 so $N=14$, $|N-10|=4$}, difficulty=3, diff_parts={K1 S2 I0 C0 F1}, diff_note={concepts: divisibility by digits | steps: enumerate by tens digit; check each; count 14}]
If the number of positive two digit integers which are divisible by both of its digit be $N$ then the value of $|N-10|$ is
\end{bankq}
\begin{bankq}[id=TOAS01-P2-MATH-S2-Q07, source={TOAS - Test 01 (30 Jan 2022)}, paper=2, qno=7, type=numerical, subtopic={Combinations}, answer={309}, answer_src=solved, working={No two consecutive from 12: $\binom{13-k}{k}$ for $k=3,4,5,6$: $120+126+56+7=309$}, difficulty=4, diff_parts={K1 S2 I1 C0 F1}, diff_note={concepts: subsets without adjacent elements | steps: formula C(13 - k, k); k = 3 to 6; evaluate; sum | I: gap method}]
Let there is a set of 12 natural numbers $A=\{481,482,\dots,492\}$. The number of subset $B$ of $A$ which contain at least 3 elements such that difference of any two elements is at least 2 is -
\end{bankq}
\begin{bankq}[id=TOAS01-P2-MATH-S2-Q08, source={TOAS - Test 01 (30 Jan 2022)}, paper=2, qno=8, type=numerical, subtopic={Distributions}, answer={4}, answer_src=solved, working={Digits $9-e_i$ with $\sum e_i=4$, last digit even so $e_6$ odd: $e_6=1$: $\binom74=35$; $e_6=3$: 5; $x=40$, $x/10=4$}, difficulty=4, diff_parts={K1 S2 I1 C0 F1}, diff_note={concepts: digit sums with constraints | steps: complement digits sum to 4; last digit even needs odd complement; cases 1 and 3; add | I: complement digits}, also_in={TNAS - Test 01 (05 Feb 2023) P1 Q10}]
Let $x$ is the number of 6 digit even numbers such that sum of their digits is 50, then $\dfrac{x}{10}$ is equal to
\end{bankq}
\begin{bankq}[id=TOAS02-P1-MATH-S3-Q01, source={TOAS - Test 02 (13 Feb 2022)}, paper=1, qno=1, type=numerical, subtopic={Counting principles}, answer={48}, answer_src=solved, working={Missing element $m$ (4 ways); onto the other three with $f(x)\ne x$: 24 maps minus 12 non-onto = 12; total 48}, difficulty=5, diff_parts={K2 S2 I1 C0 F1}, diff_note={concepts: functions with restrictions, inclusion-exclusion | steps: choose the missing element; count maps with f(x) != x on the other three elements; remove non-onto maps | I: inclusion-exclusion}, also_in={TNAS - Test 02 (07 May 2023) P1 Q7}]
Let $A=\{1,2,3,4\}$, $B=\{1,2,3,4\}$. The total number of functions $f:A\to B$ such that $f(x)\ne x\ \forall\ x\in A$ and range of $f(x)$ contains exactly three elements, is
\end{bankq}
\begin{bankq}[id=TOAS15-P1-MATH-S1-Q06, source={TOAS - Test 15 (22 Aug 2022)}, paper=1, qno=6, type=single, subtopic={Combinations}, answer={D}, answer_src=key, difficulty=1, diff_parts={K1 S1 I0 C0 F0}, diff_note={concepts: selecting shoes without a pair | steps: choose the four pairs out of five; choose one shoe from each pair}]
A rack has 5 different pairs of shoes. The number of ways in which 4 shoes can be chosen from it, such that there will be no complete pair is :
\begin{qoptions}
\opt{A}{1920}
\opt{B}{200}
\opt{C}{110}
\opt{D}{80}
\end{qoptions}
\end{bankq}
\begin{bankq}[id=TOAS15-P1-MATH-S2-Q12, source={TOAS - Test 15 (22 Aug 2022)}, paper=1, qno=12, type=multi, subtopic={Combinations}, answer={B,D}, answer_src=key, difficulty=3, diff_parts={K1 S1 I1 C0 F1}, diff_note={concepts: combinatorial interpretations of a binomial coefficient | steps: test each of the four counting statements by direct counting}]
The combinatorial coefficient $^{n-1}C_p$ denotes
\begin{qoptions}
\opt{A}{The number of ways in which $n$ things of which $p$ are alike and rest different can be arranged in a circle}
\opt{B}{The number of ways in which $p$ different things can be selected out of $n$ different thing if a particular thing is always excluded}
\opt{C}{Number of ways in which $n$ alike balls can be distributed in $p$ different boxes so that no box remains empty, and each box can hold any number of balls}
\opt{D}{The number of ways in which $(n-2)$ white balls and $p$ black balls can be arranged in a line if black balls are separated, balls are alike except for the colour}
\end{qoptions}
\end{bankq}
\begin{bankq}[id=TOAS15-P2-MATH-S1-Q02, source={TOAS - Test 15 (22 Aug 2022)}, paper=2, qno=2, type=multi, subtopic={Derangements}, answer={A,B,C,D}, answer_src=key, difficulty=4, diff_parts={K1 S2 I1 C0 F1}, diff_note={concepts: restricted permutations | steps: count permutations with no adjacent consecutive terms for small n; recurrence growth; compare Q21 and Q20}]
There are $n$-students, namely $S_1,S_2,S_3,\ldots,S_n$ and $n$-chairs placed in a row. Let $Q_n$ denotes the number of ways of seating these $n$-students such that no two students having consecutive numbers sit together. Then
\begin{qoptions}
\opt{A}{$Q_3=0$}
\opt{B}{$Q_4=2$}
\opt{C}{$Q_5=14$}
\opt{D}{$Q_{21}>19Q_{20}$}
\end{qoptions}
\end{bankq}
\begin{bankq}[id=TOAS15-P2-MATH-S3-Q06, source={TOAS - Test 15 (22 Aug 2022)}, paper=2, qno=6, type=single, subtopic={Permutations}, answer={4}, answer_src=key, difficulty=5, diff_parts={K1 S2 I2 C1 F0}, diff_note={concepts: ordered arrangements across several queues | steps: place the persons with early-service restrictions; interleave the queues; count; read k}]
$W_1$, $W_2$ and $W_3$ are three ticket windows in a cinema hall. 12 persons ($A_1,A_2,\ldots\ldots,A_{12}$) are standing for the tickets in a queue of 4, 3 and 5 number of persons infront of $W_1$, $W_2$ and $W_3$ respectively. $A_1$ and $A_2$ want to get their tickets from $W_1$ and $A_3$, $A_4$ and $A_5$ want to get their tickets from $W_3$. If the number of ways in which all are getting their tickets such that $A_1$, $A_2$, $A_3$, $A_4$ and $A_5$ are getting their tickets as early as possible is $k(7!)$, then find the value of $\left[\dfrac k3\right]$.\par
(Assuming each person is getting exactly one ticket at a time).
[Note: $[y]$ denotes greatest integer less than or equal to $y$.]
\end{bankq}
\begin{bankq}[id=TOAS17-P2-MATH-S2-Q04, source={TOAS - Test 17 (18 Sep 2022)}, paper=2, qno=4, type=numerical, subtopic={Permutations}, answer={44}, answer_src=key, difficulty=5, diff_parts={K1 S2 I1 C1 F1}, diff_note={concepts: counting invertible matrices with distinct entries | steps: total arrangements; singular cases without zero; subtract}]
Let $[a_{ij}]$ is a square matrix of order 2 where $a_{ij}\in\{0,1,2,3,4,6\}$ then number of matrices $A$ such that $A^{-1}$ exists is $\lambda$ then $\lambda-300$ is (given that all elements of matrix $A$ are distinct)
\end{bankq}
\begin{bankq}[id=TOAS17-P2-MATH-S2-Q09, source={TOAS - Test 17 (18 Sep 2022)}, paper=2, qno=9, type=numerical, subtopic={Counting principles}, answer={4}, answer_src=key, difficulty=4, diff_parts={K1 S2 I1 C0 F1}, diff_note={concepts: last digit of a power | steps: exponent modulo 4; power of three; add one}]
The last digit of $3^{81^n}+1$, $n\in N$, is
\end{bankq}
\begin{bankq}[id=TOAS19-P1-MATH-S1-Q04, source={TOAS - Test 19 (05 Feb 2023)}, paper=1, qno=4, type=multi, subtopic={Counting principles}, answer={A,B}, answer_src=key, difficulty=5, diff_parts={K1 S2 I1 C1 F1}, diff_note={concepts: lattice path counting with steps of 1 or 2 | steps: recursion for the number of ways; compute the grid; compare}, also_in={TNAS - Test 01 (05 Feb 2023) P2 Q4}]
In a game a person $P$ at point $(x,y)$ can jump to any one of the following points relatively: $(x+1,y)$, $(x+2,y)$, $(x,y+1)$ or $(x,y+2)$. If $P$ is at $(0,0)$, then he can reach
\begin{qoptions}
\opt{A}{$(2,2)$ in 14 ways}
\opt{B}{$(2,3)$ in 32 ways}
\opt{C}{$(2,2)$ in 9 ways}
\opt{D}{$(2,3)$ in 26 ways}
\end{qoptions}
\end{bankq}

\section*{TOPS papers}
\begin{bankq}[id=TOPS05-P1-MATH-S3-Q02, source={TOPS - Test 05 (14 Jan 2023)}, paper=1, qno=2, type=numerical, subtopic={Permutations}, answer={720.00}, answer_src=key, difficulty=6, diff_parts={K1 S2 I2 C1 F1}, diff_note={concepts: permutations with repetition | steps: choose the repeated digit; positions; remaining distinct digits; count}]
The number of 5-digit nutural numbers beginning with 9 that has exactly three identical digits (other digits are distinct) is
\end{bankq}
\begin{bankq}[id=TOPS05-P1-MATH-S3-Q05, source={TOPS - Test 05 (14 Jan 2023)}, paper=1, qno=5, type=numerical, subtopic={Permutations}, answer={1.5}, answer_src=solved, working={Last digit 5: $\dfrac{7!}{2!2!2!}-\dfrac{6!}{2!2!}=450$. Last digit 3: $\dfrac{7!}{2!2!}-\dfrac{6!}{2!}=900$. Total $1350$, divided by 900 gives $1.5$}, difficulty=7, diff_parts={K1 S3 I2 C1 F1}, diff_note={concepts: permutations with repetition | steps: odd last digit; arrangements; exclude zero first; both cases; divide}]
Let $\lambda$ be number of all possible 8 digit odd numbers formed by using only digits 0, 0, 2, 2, 3, 3, 4, 5, then $\dfrac\lambda{900}$ is
\end{bankq}
\begin{bankq}[id=TOPS05-P1-MATH-S3-Q08, source={TOPS - Test 05 (14 Jan 2023)}, paper=1, qno=8, type=numerical, subtopic={Combinations}, answer={96.00}, answer_src=key, difficulty=6, diff_parts={K2 S2 I1 C1 F1}, diff_note={concepts: polygon diagonals, counting | steps: n from the diagonals; triangles with one side; count}]
A regular polygon of $n$ sides has 54 diagonals. The number of triangles formed with exactly one side of the triangle coinciding with the side of polygon, is
\end{bankq}
\begin{bankq}[id=TOPS05-P2-MATH-S1-Q04, source={TOPS - Test 05 (14 Jan 2023)}, paper=2, qno=4, type=single, subtopic={Permutations}, answer={C}, answer_src=key, difficulty=5, diff_parts={K1 S2 I2 C1 F0}, diff_note={concepts: permutations with restrictions | steps: count with no restriction; choose the allowed letters; arrange; total}]
The number of ways of arranging letters of the word COMPUTER such that none of the first three letters are $C$, $O$ or $M$, is
\begin{qoptions}
\opt{A}{7100}
\opt{B}{7000}
\opt{C}{7200}
\opt{D}{7120}
\end{qoptions}
\end{bankq}
\begin{bankq}[id=TOPS05-P2-MATH-S3-Q02, source={TOPS - Test 05 (14 Jan 2023)}, paper=2, qno=2, type=numerical, subtopic={Permutations}, answer={7}, answer_src=key, difficulty=7, diff_parts={K2 S2 I2 C1 F1}, diff_note={concepts: permutations with repetition, no adjacent | steps: total arrangements; restrictions on equal letters; inclusion exclusion; factorise}]
If number of arrangements of letters of the word ``DHARAMSHALA'' taken all at a time so that no two alike letters appear together is $\left(4^a.5^b.6^c.7^d\right)$, (where $a,b,c,d\in N$), then $a+b+c+d$ is equal to
\end{bankq}

\section*{TMAS papers}
\begin{bankq}[id=TMAS04-P1-MATH-S3-Q01, source={TMAS - Test 04 (19 Nov 2023)}, paper=1, qno=1, type=numerical, subtopic={Counting principles}, answer={48.00}, answer_src=key, difficulty=6, diff_parts={K1 S2 I2 C1 F1}, diff_note={concepts: counting functions with a restriction | steps: choose the missing value; assign; subtract fixed points; total}]
Let $A=\{1,2,3,4\}$, $B=\{1,2,3,4\}$. The total number of functions $f:A\to B$ such that $f(x)\neq x\ \forall\ x\in A$ and range of $f(x)$ contains exactly three elements, is
\end{bankq}

\section*{TNAS papers}
\begin{bankq}[id=TNAS01-P1-MATH-S1-Q05, source={TNAS - Test 01 (05 Feb 2023)}, paper=1, qno=5, type=single, subtopic={Counting principles}, answer={A}, answer_src=key, difficulty=4, diff_parts={K1 S2 I1 C1 F0}, diff_note={concepts: arithmetic progression selection | steps: common difference; counts; sum}]
The number of ways in which three distinct numbers in A.P. can be selected from $1,2,3,\ldots\ldots,24$ is -
\begin{qoptions}
\opt{A}{$132$}
\opt{B}{$572$}
\opt{C}{$264$}
\opt{D}{$150$}
\end{qoptions}
\end{bankq}
\begin{bankq}[id=TNAS01-P1-MATH-S1-Q08, source={TNAS - Test 01 (05 Feb 2023)}, paper=1, qno=8, type=single, subtopic={Permutations}, answer={C}, answer_src=key, difficulty=4, diff_parts={K1 S2 I1 C1 F0}, diff_note={concepts: arrangement with conditions | steps: treat the pair as one; separate the others; count}]
The number of ways in which the letters of the word `LONDON' can be rearranged if the two `O' s are together but the two `N's are seperated is
\begin{qoptions}
\opt{A}{$72$}
\opt{B}{$71$}
\opt{C}{$36$}
\opt{D}{$35$}
\end{qoptions}
\end{bankq}
\begin{bankq}[id=TNAS01-P1-MATH-S1-Q09, source={TNAS - Test 01 (05 Feb 2023)}, paper=1, qno=9, type=single, subtopic={Counting principles}, answer={B}, answer_src=key, difficulty=3, diff_parts={K1 S1 I1 C1 F0}, diff_note={concepts: highest power in a factorial | steps: powers of 2 and 3; minimum}]
If $(201)!$ is divided by $24^k$, then the largest value of $k$ is
\begin{qoptions}
\opt{A}{$98$}
\opt{B}{$65$}
\opt{C}{$49$}
\opt{D}{$66$}
\end{qoptions}
\end{bankq}
\begin{bankq}[id=TNAS01-P1-MATH-S1-Q13, source={TNAS - Test 01 (05 Feb 2023)}, paper=1, qno=13, type=single, subtopic={Combinations}, answer={C}, answer_src=key, difficulty=3, diff_parts={K1 S1 I1 C1 F0}, diff_note={concepts: selection with repetition | steps: cases of repeated letters; count}]
The number of ways of selection of $4$ letters of the word `PERFECT ' is
\begin{qoptions}
\opt{A}{$15$}
\opt{B}{$21$}
\opt{C}{$25$}
\opt{D}{$31$}
\end{qoptions}
\end{bankq}
\begin{bankq}[id=TNAS01-P1-MATH-S1-Q14, source={TNAS - Test 01 (05 Feb 2023)}, paper=1, qno=14, type=single, subtopic={Counting principles}, answer={A}, answer_src=key, difficulty=3, diff_parts={K1 S1 I1 C1 F0}, diff_note={concepts: lattice points in a circle | steps: count by value of x; sum}]
The number of points $(x,y)$ having integral coordinates satisfying $x^2+y^2<25$ is -
\begin{qoptions}
\opt{A}{$69$}
\opt{B}{$80$}
\opt{C}{$81$}
\opt{D}{$77$}
\end{qoptions}
\end{bankq}
\begin{bankq}[id=TNAS01-P1-MATH-S1-Q17, source={TNAS - Test 01 (05 Feb 2023)}, paper=1, qno=17, type=single, subtopic={Permutations}, answer={A}, answer_src=key, difficulty=3, diff_parts={K1 S1 I1 C1 F0}, diff_note={concepts: sum of formed numbers | steps: contribution of each place; total}]
The sum of all $4$ digit numbers that can be formed by using the digits $2,4,6,8$. (repetition of digits not being allowed) is -
\begin{qoptions}
\opt{A}{$133320$}
\opt{B}{$123330$}
\opt{C}{$113230$}
\opt{D}{$323430$}
\end{qoptions}
\end{bankq}
\begin{bankq}[id=TNAS01-P1-MATH-S1-Q20, source={TNAS - Test 01 (05 Feb 2023)}, paper=1, qno=20, type=single, subtopic={Counting principles}, answer={B}, answer_src=key, difficulty=5, diff_parts={K1 S2 I2 C1 F0}, diff_note={concepts: sum of divisors | steps: odd divisors; divisible by five not seven; sum formula}]
Consider the number $N=2^4\,3^4\,5^2\,7^2\,11^2$, then the sum of all the odd divisors which are divisible by `5' but not by `7' is $11^2(133\lambda)$ then $\lambda$ is equal to
\begin{qoptions}
\opt{A}{$10$}
\opt{B}{$30$}
\opt{C}{$50$}
\opt{D}{$330$}
\end{qoptions}
\end{bankq}

\section*{Nurture (General) papers}
\begin{bankq}[id=NUR01-P2-MATH-S3-Q01, source={Nurture (General) - Test 01 (12 Jun 2021)}, paper=2, qno=1, type=single, subtopic={Combinations}, answer={C}, answer_src=solved, working={k is present and chosen freely from smaller elements: 2\^{}(k-1)}, difficulty=2, diff_parts={K1 S1 I1 C0 F0}, diff_note={concepts: subsets with a given largest element | steps: fix the largest; choose smaller elements}]
For any given $n$, number of subsets of $A$ having $k(k<n)$ as largest element, is
\begin{qoptions}
\opt{A}{$2^k$}
\opt{B}{$n-k$}
\opt{C}{$2^{k-1}$}
\opt{D}{None of these}
\end{qoptions}
\end{bankq}
\begin{bankq}[id=NUR01-P2-MATH-S3-Q02, source={Nurture (General) - Test 01 (12 Jun 2021)}, paper=2, qno=2, type=single, subtopic={Combinations}, answer={C}, answer_src=solved, working={Total minus subsets with only even elements: 2\^{}n-2\^{}[n/2]}, difficulty=3, diff_parts={K1 S2 I1 C0 F0}, diff_note={concepts: subsets with an odd element | steps: complement; even elements only; count}]
The number of subsets of $A$ having at least one odd integer is [where $[.]$ represents greatest integer function]
\begin{qoptions}
\opt{A}{$2^{\left[\frac n2\right]}$}
\opt{B}{$2^{\left[\frac{n+1}2\right]}$}
\opt{C}{$2^n-2^{\left[\frac n2\right]}$}
\opt{D}{$2^{\left[\frac n2\right]+1}$}
\end{qoptions}
\end{bankq}
\begin{bankq}[id=NUR04-P1-MATH-S1-Q03, source={Nurture (General) - Test 04 (11 Feb 2024)}, paper=1, qno=3, type=single, subtopic={Distributions}, answer={C}, answer_src=solved, working={Sums (7,11): 15x10; (11,7): 45x6; total 150+270=420}, difficulty=6, diff_parts={K2 S2 I2 C1 F0}, diff_note={concepts: positive integral solutions of a product, distributions | steps: factor pairs; stars and bars for each; add}]
Number of positive integral solutions satisfying the equation $(x_1+x_2+x_3)(y_1+y_2)=77$, is:
\begin{qoptions}
\opt{A}{150}
\opt{B}{270}
\opt{C}{420}
\opt{D}{1024}
\end{qoptions}
\end{bankq}
\begin{bankq}[id=NUR04-P1-MATH-S2-Q10, source={Nurture (General) - Test 04 (11 Feb 2024)}, paper=1, qno=10, type=multi, subtopic={Distributions}, answer={A,C,D}, answer_src=solved, working={6!/(2!2!2!)=90: apples/mangoes 6x15, the multinomial coefficient and the equal distribution give 90; bundles give 15}, difficulty=7, diff_parts={K2 S2 I2 C1 F1}, diff_note={concepts: counting identity with repeated letters, distributions | steps: arrangements; each alternative; compare}]
Number of ways in which the letters of the word `B U L B U L' can be arranged in a line in any order is also equal to the
\begin{qoptions}
\opt{A}{Number of ways in which 2 alike Apples and 4 alike Mangoes can be distributed in 3 children so that each child receives any number of fruits.}
\opt{B}{Number of ways in which 6 different books can be tied up into 3 bundles, if each bundle is to have equal number of books}
\opt{C}{Coefficient of $x^2y^2z^2$ in the expansion of $(x+y+z)^6$}
\opt{D}{Number of ways in which 6 different prizes can be distributed equally in three children.}
\end{qoptions}
\end{bankq}
\end{document}
